
% =============================================================================
% Figure 02 — Domains of Invariance
% Book I — Ontology of Viability
% Chapter 02 — Invariance
%
% Purpose:
%   To illustrate the bounded and conditional scope within which
%   an invariant may obtain under specified transformations.
%
% Conceptual Policy:
%   1. Domains of invariance are analytical scopes, not ontological
%      realms or independent components of existence.
%   2. Invariance is always relative to a specified respect.
%   3. Different domains may overlap.
%   4. Validity within one domain does not imply validity in another.
%   5. No invariant is presumed to hold across all transformations.
%   6. The diagram does not establish a hierarchy among domains.
% =============================================================================

\begin{figure}[htbp]
    \centering

    \begin{tikzpicture}[
        font=\small,
        domain/.style={
            draw,
            rounded corners=4pt,
            thick,
            minimum width=3.35cm,
            minimum height=1.35cm,
            align=center
        },
        central/.style={
            draw,
            rounded corners=4pt,
            thick,
            minimum width=4.0cm,
            minimum height=1.25cm,
            align=center,
            fill=gray!10
        },
        scope/.style={
            draw,
            dashed,
            rounded corners=5pt,
            thick,
            inner sep=0.55cm
        },
        relation/.style={
            thick,
            ->,
            >=stealth
        },
        possible/.style={
            thick,
            dashed,
            ->,
            >=stealth
        }
    ]

    % -------------------------------------------------------------------------
    % Central analytical claim
    % -------------------------------------------------------------------------

    \node[central] (claim) {
        \textbf{Invariance Claim}\\
        A specified feature is preserved
    };

    % -------------------------------------------------------------------------
    % Domains
    % -------------------------------------------------------------------------

    \node[domain, fill=blue!6, above left=1.55cm and 2.25cm of claim]
        (structural) {
        \textbf{Structural Domain}\\
        Preservation of organization
    };

    \node[domain, fill=green!7, above right=1.55cm and 2.25cm of claim]
        (relational) {
        \textbf{Relational Domain}\\
        Preservation of relations
    };

    \node[domain, fill=orange!9, below left=1.55cm and 2.25cm of claim]
        (determinational) {
        \textbf{Determinational Domain}\\
        Preservation of determinations
    };

    \node[domain, fill=red!6, below right=1.55cm and 2.25cm of claim]
        (compositional) {
        \textbf{Compositional Domain}\\
        Preservation of composition
    };

    % -------------------------------------------------------------------------
    % Relations between claim and domains
    % -------------------------------------------------------------------------

    \draw[possible] (claim) -- (structural);
    \draw[possible] (claim) -- (relational);
    \draw[possible] (claim) -- (determinational);
    \draw[possible] (claim) -- (compositional);

    % -------------------------------------------------------------------------
    % Domain of validity
    % -------------------------------------------------------------------------

    \node[
        draw,
        rounded corners=4pt,
        thick,
        fill=gray!5,
        minimum width=10.2cm,
        minimum height=0.85cm,
        align=center,
        below=1.55cm of claim
    ] (validity) {
        \textbf{Specified Domain of Validity}\\
        Transformation, conditions, and respect of preservation
    };

    \draw[relation] (claim) -- (validity);

    % -------------------------------------------------------------------------
    % Boundary note
    % -------------------------------------------------------------------------

    \node[
        below=0.65cm of validity,
        text width=11cm,
        align=center,
        font=\footnotesize
    ] (note) {
        Invariance established within a specified domain
        cannot be generalized automatically to other domains
        or to unrestricted transformations.
    };

    \end{tikzpicture}

    \caption{Analytical Domains of Invariance.
    An invariance claim may concern different kinds of preserved
    determinations, relations, or organizations. Its validity
    remains bounded by the specified transformation, conditions,
    and respect of preservation.}

    \label{fig:domains-of-invariance}

\end{figure}
